%! Characteristics of Polynomial Functions — Unit 4, Lesson 4.0 — Student Packet Cover Sheet
\documentclass[10pt]{article}
\usepackage{algebra2-article}
\usepackage{algebra2-boxes}
\usepackage{ltablex}
\keepXColumns

\begin{document}

% Full-bleed forest banner
\begin{tikzpicture}[remember picture, overlay]
  \fill[forest] (current page.north west)
    rectangle ([yshift=-0.9in]current page.north east);
\end{tikzpicture}
\vspace{-0.8in}

\begin{center}
  {\color{white}\textbf{\LARGE Algebra 2: Shepherd}}\\[4pt]
  {\color{forestmid}\large\textbf{Unit 4 \quad Polynomial Functions}}\\[6pt]
  {\color{white}\textbf{\large Lesson 4.0 \quad Characteristics of Polynomial Functions}}
\end{center}

\vspace{0.22in}
\namedateperiod
\vspace{0.18in}

\begin{learningtargetbox}
\small
\begin{enumerate}[leftmargin=1.5em, itemsep=5pt]
  \item \ldots{} identify a polynomial's \textbf{degree} and \textbf{leading coefficient}, and use them to describe its \textbf{end behavior}.
  \item \ldots{} read a polynomial's \textbf{turning points}, its \textbf{relative} and \textbf{absolute} maxima/minima, its \textbf{domain}, \textbf{range}, \textbf{intercepts}, and \textbf{zeros}, and use each zero's \textbf{multiplicity} to tell whether the graph \textbf{crosses} or \textbf{touches} the $x$-axis.
  \item \ldots{} classify a polynomial's \textbf{symmetry} as \textbf{even} ($y$-axis), \textbf{odd} (origin), or neither, and justify it.
\end{enumerate}
\end{learningtargetbox}

\vspace{0.14in}

\begin{tocbox}
\small
\renewcommand{\arraystretch}{1.5}
\begin{tabularx}{\linewidth}{@{} c l X r @{}}
\rowcolor{forestbg}
\textbf{\#} & \textbf{Component} & \textbf{Description} & \textbf{Score} \\
\midrule
1 & Warm-Up        & Spiral review: powers of negatives, evaluating a cubic, zeros from factors & \blank{1.2cm} \\
2 & Guided Notes   & Reading every characteristic off a polynomial graph                        & \blank{1.2cm} \\
3 & Group Activity & \emph{Reading Polynomial Graphs} --- feature lists and a volume model, by tier & \blank{1.2cm} \\
4 & Exit Ticket    & Quick check: end behavior, extrema, multiplicity, and an SOL-style item     & \blank{1.2cm} \\
5 & Homework       & Feature lists from factored forms and graphs, plus a modeling extension     & \blank{1.2cm} \\
\midrule
  & \multicolumn{2}{r}{\textbf{Total}} & \blank{1.2cm} \\
\end{tabularx}
\end{tocbox}

\vspace{0.10in}

\begin{remindbox}
\small A \textbf{polynomial function} $f(x)=a_nx^n+\dots+a_1x+a_0$ graphs as a \textbf{smooth,
continuous curve} that can rise and fall several times. Its whole ``personality'' can be read off
that curve: its \textbf{end behavior} (set by the \emph{degree} --- even or odd --- and the sign of
the \emph{leading coefficient}), its \textbf{turning points} (at most $n-1$ of them) and the
\textbf{relative} or \textbf{absolute} maxima/minima there, its \textbf{$y$-intercept} and its
\textbf{zeros} (where each zero's \textbf{multiplicity} decides whether the graph \emph{crosses} or
\emph{touches} the $x$-axis), its \textbf{domain} and \textbf{range}, where it is \textbf{increasing}
or \textbf{decreasing} and \textbf{positive} or \textbf{negative}, and its \textbf{symmetry} (even
about the $y$-axis, odd about the origin, or neither). Every one of these is a picture you can point
to.
\end{remindbox}

\end{document}
